\subsection{Vacuum Fluctuations and the Casimir Effect}
  \label{subsec:vacuum-fluctuations-and-the-casimir-effect}

  Vacuum fluctuations arise from the non-injectivity of the projection~$\Pi$ in
  regimes where no stable localized excitations are present.
  In such regimes, a wide multiplicity of distinct relational configurations of~$\chi$
  project onto the same effective spacetime description.
  Any two such co-projecting configurations $\omega_1, \omega_2 \in \Pi^{-1}(o)$
  share a common projected part---which yields the stable observable~$o$---and carry
  residual structural differences that are not erased by~$\Pi$.
  The projected manifestation of these residual differences constitutes the observable
  fluctuation.
  No intrinsic indeterminacy of the pre-geometric substrate is required: $\chi$ does not
  vary at the substrate level.
  Fluctuations are therefore an exclusive feature of the projected description, and their
  magnitude is controlled by the degree of non-injectivity of~$\Pi$, quantified by the
  effective projective entropy~$S_\Pi$.

  The spatiotemporal character of fluctuations---their variation in both space and
  time---follows directly from the same mechanism.
  Spatial variations arise from residual relational differences between co-projecting
  configurations that are localized in distinct regions of the effective description.
  Temporal variations arise because co-projecting configurations may differ in their
  position along the relaxation-ordering parameter~$\lambda$: such ordering differences,
  once projected, are precisely what the effective description identifies as differences in
  operational time~$\tau$.
  The time-dependence of vacuum fluctuations is therefore not a signature of any
  infra-physical dynamics of~$\chi$, but the effective translation of relaxation-ordering
  differences between co-projecting configurations.

  A natural question arises at this point: why do projection residuals \emph{appear random}
  rather than repeating identically under successive projections?
  If the substrate~$\chi$ were stationary and the projection~$\Pi$ were applied repeatedly to
  the same fibre configuration, the residuals would indeed repeat identically.
  The appearance of randomness is therefore not a property of~$\chi$ itself, but a consequence
  of the fact that the effective projection is not applied at a fixed relaxation stage.
  More precisely, the projection must be understood as implicitly $\lambda$-dependent:
  one effectively probes a family of maps~$\Pi_\lambda$.
  At each effective step~$U_n$, one probes not a fixed fibre $\Pi^{-1}(o)$, but a slice of the
  co-projecting set determined by the current value of the relaxation-ordering
  parameter~$\lambda$.
  Between two successive steps $U_n$ and $U_{n+1}$, $\lambda$ has advanced, and the set of
  substrate configurations that co-project onto the same observable~$o$ is therefore not
  identical.
  The advance of~$\lambda$ is monotonic but remains unobservable at the effective level,
  since $\lambda$-ordering is precisely what the effective description identifies as time.
  Each reprojection thus samples a \emph{different slice} of an evolving fibre, and the
  residual differences between these slices are registered as fluctuations in the effective
  description.
  In this sense, the apparent stochasticity of fluctuations is the observable signature of the
  irreversible advance of~$\lambda$, rather than an intrinsic indeterminacy of~$\chi$.
  Equivalently, if relaxation could be suspended while reprojections continued, the same
  residuals would repeat identically.
  Fluctuations therefore arise because the projection cannot be iterated at fixed~$\lambda$:
  time, understood as the projected image of relaxation ordering, is the minimal and
  irreducible source of the observed noise.

  \begin{remark}[Fluctuations as $\lambda$-induced fibre drift]
  The mechanism identified above admits a natural geometric reading.
  The family $\{\Pi_\lambda\}_{\lambda \in \Lambda}$ defines a $\lambda$-parametrised bundle
  of co-projecting fibres over the observable space~$\mathcal{O}$.
  Fluctuations correspond to the transverse variation of successive fibre slices along the
  relaxation flow, i.e.\ to the drift of~$\Pi_\lambda^{-1}(o)$ as $\lambda$ advances.
  A rigorous formulation of this drift---as a connection on the fibre bundle, or as a
  Lie-derivative along the relaxation vector field---is left as an open structural direction.
  \end{remark}

  Two structural predictions follow from this picture.
  First, in a regime where the effective projection becomes \emph{locally} injective,
  the local projective entropy $S_\Pi$ is suppressed and quantum fluctuations are
  correspondingly reduced.
  Since the ordering differences between co-projecting configurations are precisely
  what generate distinguishable successive effective states, a locally injective regime
  implies $U_{n+1} = U_n$ in that region: the projected description becomes locally
  stationary and operational time ceases to advance there.
  This is the projective reading of decoherence and of the classical limit: it does not
  require environmental averaging as a primitive assumption, but follows from the
  local reduction of the co-projecting fiber.
  Such locally injective regimes are admissible; they correspond to what is ordinarily
  described as the quasi-classical regime.
  A locally injective regime is one in which the restriction of~$\Pi$ to the accessible
  configuration neighborhood admits a single dominant element of~$\Pi^{-1}(o)$:
  distinct substrate configurations within that neighborhood produce distinguishable
  projected states, even though the global projection remains non-injective.

  Second, a \emph{globally} injective projection---one in which $|\Pi^{-1}(o)| = 1$
  for all observables simultaneously---would eliminate all projective entropy and
  terminate temporal ordering everywhere.
  This is structurally excluded on two independent grounds.
  The configuration space $\Omega_\chi$ of the substrate is structurally richer than the
  observable space $O$: the projection $\Pi$ is an irreducible coarse-graining, so
  $|\Omega_\chi| \gg |O|$ and a globally injective $\Pi$ is mathematically incompatible
  with the compression implicit in any effective description.
  Furthermore, a globally stationary projection would imply $U_{n+1} = U_n$ for all
  effective observers simultaneously, which is the cosmological counterpart of absolute
  zero and is excluded by the same argument developed in Section~5.6.
  The persistence of a minimal global non-injectivity is therefore not merely compatible
  with temporal evolution; it is a structural requirement for it.

  When material boundaries impose structural constraints on local
  projectability, certain effective descriptions become incompatible with
  the boundary conditions, reducing the set of admissible projective
  configurations between the boundaries.
  The Casimir effect arises from this asymmetry: a difference in the
  density of admissible effective reprojections, manifesting as a pressure
  on the confining surfaces.
  No fundamental vacuum energy density or propagating vacuum modes are
  required.
